%% kapa: exam=1 round=25 year=2014 period=2 form=A subject=accounting subject_ko=회계학 seq=27 number=67
%% source: Q-Net 감정평가사 기출문제 (cst00309) artlSeq=5101207 file=extracted/1008/5101207/2014년 제25회 감정평가사 1차시험 기출문제/2교시(A형).hwp
%% conversion: kordoc 4.24.51 -> tools/gen1.py (auto); review=proofread-sonnet; flags=none
\begin{question}{67}
20×1년 12월 1일 원화가 기능통화인 (주)서울은 해외 거래처에 US \$5,000의 상품을 판매하고 판매대금은 2개월 후인 20×2년 1월 31일에 회수하였다. 이 기간 중 US \$ 대비 원화의 환율은 아래와 같으며, 회사는 회계기준에 준거하여 외화거래 관련 회계처리를 적절하게 수행하였다.
\begin{framedtext}
\hspace*{3em}20×1년 12월 \ 1일: US \$1 = ￦1,030 \par
\hspace*{3em}20×1년 12월 31일: US \$1 = ￦1,060 \par
\hspace*{3em}20×2년 \ 1월 31일: US \$1 = ￦1,050 \par
\end{framedtext}
대금결제일인 20×2년 1월 31일에 (주)서울이 인식할 외환차익 혹은 외환차손은?
\begin{choices}
  \item 외환차손 ￦50,000
  \item 외환차손 ￦100,000
  \item 외환차익 ￦100,000
  \item 외환차익 ￦150,000
  \item 외환차손 ￦150,000
\end{choices}
\end{question}
\answer{1}
